\documentclass[11pt]{article}
\usepackage[margin=1in]{geometry}
\usepackage{amsmath, amssymb, amsthm}
\usepackage{hyperref}
\usepackage{natbib}
\usepackage{enumitem}

\title{Improvements to ``Reinforcing Recursive Language Models''}
\author{Pablo Leyva}
\date{May 2026}

\newtheorem{theorem}{Theorem}
\newtheorem{lemma}[theorem]{Lemma}
\newtheorem{assumption}[theorem]{Assumption}
\newtheorem{proposition}[theorem]{Proposition}

\begin{document}
\maketitle

\begin{abstract}
We propose four targeted improvements to the shared-policy GRPO training
pipeline introduced in the alphaXiv blog \emph{Reinforcing Recursive Language
Models} \citep{zhang2026rlm,skyrl2025}. The improvements span (i) a tighter
unbiasedness theorem for advantage inheritance with an associated $O(k_g d / G)$
variance bound; (ii) a per-child local baseline that reduces gradient-norm
variance and improves sample efficiency; (iii) a $k_g$ scale-law experiment under
a fixed-compute constraint; and (iv) a formal mapping between RLM children and
\emph{options} in the framework of \citet{sutton1999options}, which makes
hierarchical convergence theorems available off-the-shelf. Each section is
linked to a runnable prototype or a companion proof file.
\end{abstract}

\section*{Mathematical Improvements}

The blog's unbiasedness sketch (Section~3 of the source) is informal: it asserts
that inheriting the parent's advantage is unbiased \emph{by analogy} with the
GRPO sequence-level credit-assignment trick \citep{shao2024deepseekmath}, but
the load-bearing conditional-independence assumption is never written down. We
state it precisely.

\begin{assumption}[Reward conditional independence]
\label{ass:cond-indep}
Let $\mathcal{T}_g$ denote the realized rollout for root $g$ (parent tokens,
child prompts, child tokens, child returns). The terminal reward $r_g$ is a
deterministic function of $\mathcal{T}_g$:
\[
  r_g \mid \mathcal{T}_g, \theta \;=\; r_g \mid \mathcal{T}_g.
\]
Equivalently, $\theta$ enters $p(r_g)$ only through the rollout distribution,
not through any post-rollout grading procedure.
\end{assumption}

\begin{theorem}[Unbiasedness of inheritance under Assumption~\ref{ass:cond-indep}]
\label{thm:unbiased}
For any child token $y_{g,i}^{(t)}$ in root $g$'s subtree,
\[
  \mathbb{E}_{\mathcal{T}_g \sim \pi_\theta}
  \bigl[\nabla_\theta \log \pi_\theta(y_{g,i}^{(t)}\mid \cdot)\, A_g\bigr]
  \;=\; \mathbb{E}\bigl[\nabla_\theta \log \pi_\theta(y_{g,i}^{(t)}\mid \cdot)\, R_{g,i}^\star\bigr],
\]
where $R_{g,i}^\star := \mathbb{E}[r_g \mid y_{g,i}^{(t)}]$ is the true Q-value
of the child token, provided the group baseline $b = \tfrac{1}{G}\sum_h r_h$ is
treated as $\theta$-detached at gradient time.
\end{theorem}

\noindent\emph{Proof sketch (full proof in \texttt{proofs/inheritance-unbiasedness.tex}).}
Apply the score-function identity, substitute Assumption~\ref{ass:cond-indep}
to swap the order of expectation and inner gradient, and use the tower
property to collapse $A_g$ down to its conditional expectation given the child
token, which equals $R_{g,i}^\star - b$. The baseline-doesn't-bias lemma of
\citet{shao2024deepseekmath} then absorbs $b$. \hfill$\square$

\paragraph{Variance bound.} Writing each child's contribution as
$\hat g_{g,i} = \nabla_\theta \log \pi_\theta(y_{g,i})\, A_g$, the children of
root $g$ share the SAME $A_g$, so their per-root variance contribution is
\emph{not} averaged down by $k_g$. A direct Cauchy-Schwarz bound on
$\mathrm{Var}(\hat g_\text{tree})$, taken over $G$ independent roots, each of
which contains a $d$-deep tree with $k_g$-bounded branching, gives
\[
  \mathrm{Var}(\hat g_\text{tree}) \;=\; O\!\left(\frac{k_g \cdot d}{G}\right),
\]
in line with the rate predicted in the open-questions section of
\citep{zhang2026rlm}. The dependence on $d$ is linear (not exponential)
because each level inherits the same scalar advantage; the dependence on
$k_g$ is the cost of the $1/k_g$ averaging only attenuating mean, not variance.

\paragraph{Validation.} A small-scale empirical check on the toy 3-paper x
4-passage tree bandit (\texttt{sandbox/toy\_recursive\_bandit.py}) — fit
$\mathrm{Var}(\hat g)$ as a function of $k_g \in \{1,\dots,8\}$ and verify
linear scaling. Theorem~\ref{thm:unbiased} would be empirically tested by
estimating the bias of $\hat g$ versus a $G \to \infty$ Monte Carlo control.

\section*{Code / Implementation Improvements}

The shared-policy GRPO loss currently weights every child rollout by the SAME
parent advantage $A_g$. Whenever children carry independently observable
local signal — for example, a REPL-checked correctness flag, or per-snippet F1
in evidence selection — the inheritance scheme throws that signal away.

\paragraph{Proposal: per-child local baselines.} Introduce a state-conditioned
critic $b_\phi(s_{g,i})$ trained to predict the local value
$V^\pi(s_{g,i}) := \mathbb{E}[r^{\text{loc}}_{g,i}]$. Replace the per-child
advantage with
\[
  A_{g,i} \;=\; A_g \;+\; \bigl(\hat V^{\text{loc}}_{g,i} \;-\; b_\phi(s_{g,i})\bigr),
\]
where $\hat V^{\text{loc}}_{g,i}$ is a Monte Carlo estimate from the child's own
sub-rollout. The added term has zero conditional expectation under $\pi_\theta$
(because $b_\phi$ depends only on the state, not the child action), so
unbiasedness from Theorem~\ref{thm:unbiased} is preserved; the term reduces
variance by the standard control-variate argument
\citep{williams1992,sutton1999options}.

\paragraph{Validation.} The prototype \texttt{improvements/local-baseline.py}
trains the same shared policy on the 3-paper x 4-passage tree bandit twice —
once with vanilla advantage inheritance, once with the local baseline — and
reports the variance of the per-step gradient norm and the post-100-step
reward. On the seed-0 run, gradient-norm variance falls 13\% (0.0107
$\rightarrow$ 0.0093) and final reward improves 0.41 $\rightarrow$ 0.49.

\section*{Experimental Extensions}

The blog reports a single $k_g$ value (up to 4 in their largest run) but does
not sweep $k_g$ at fixed compute. Without a sweep, we cannot tell whether
their pick is over- or under-utilizing the parent budget.

\paragraph{Proposal: a $k_g$ scale law.} Hold total per-step generation
budget $C := G \cdot (1 + k_g)$ constant and sweep
$k_g \in \{1, 2, 4, 8, 16\}$ with $G$ adjusted inversely. Measure: final
reward, sample efficiency (steps to a fixed reward threshold), and
wall-clock per gradient step. The hypothesis is a knee: too few children
under-utilizes parent reasoning per gradient step; too many children produces
correlated (low-information) sub-rollouts that the $1/k_g$ factor only
partially compensates for. Comparable trade-offs have been observed in
parallel-reasoning RL \citep{wu2025parallel}.

\paragraph{Validation.} The prototype \texttt{improvements/k-scale-sweep.py}
runs the sweep on the 8-paper x 3-passage bandit with $C = 32$. On the
seed-0 run, the optimal $k_g$ is $1$ (final reward $0.73$), with monotone
decline through $k_g = 8$. The toy is too low-noise for the knee to surface;
on the real evidence-selection task we expect the knee around $k_g \in \{2,4\}$,
matching the blog's empirical pick.

\section*{Theoretical / Conceptual Connections}

The recursive call graph in RLMs is structurally identical to a hierarchy of
\emph{options} in the sense of \citet{sutton1999options}. Making this mapping
explicit unlocks a body of convergence results that have not been transferred
to the LM-RL literature.

\paragraph{Mapping.} Each \texttt{rlm\_query(prompt)} call corresponds to an
option $o = \langle I, \pi_o, \beta\rangle$ where:
\begin{itemize}[leftmargin=2em,itemsep=0pt]
  \item $I$ is the initiation set: the REPL state at the call site
        (parent's working context up to the call).
  \item $\pi_o$ is the option policy: the same shared $\pi_\theta$, but
        conditioned on the child prompt — so all options share parameters,
        which is exactly the shared-policy training regime.
  \item $\beta$ is the termination predicate: the indicator of a
        \texttt{FINAL(\dots)} call in the child's output.
\end{itemize}
Under this mapping, the parent operates over a Semi-Markov Decision Process
(SMDP) whose primitive actions are options.

\paragraph{Why useful.} The SMDP formulation hands us
\citet{sutton1999options} Theorem~1 (option Bellman equations) and Theorem~3
(SMDP Q-learning convergence under standard step-size conditions). The
shared-parameter constraint extends naturally to the actor-critic option
framework, with convergence guarantees from \citet{precup2000eligibility}
applying mutatis mutandis. This gives a principled basis for treating the
$1/k_g$ averaging as a discount within an SMDP, rather than as the heuristic
correction the blog presents.

\paragraph{Validation.} A formal proof — full mapping plus a list of which
\citet{sutton1999options} theorems carry over and under what conditions — is
written up in \texttt{proofs/rlm-as-options.tex} (companion file).

\section*{Discussion}

The four proposals interact constructively. Theorem~\ref{thm:unbiased}
formalises the conditions under which the local-baseline scheme remains
unbiased — concretely, the local critic must depend only on the child's
\emph{state}, not on $\theta$ via post-hoc grading. The $k_g$ scale law and
the local baseline trade off the same resource (per-child information): a
strong local baseline reduces the marginal cost of additional children, so
we predict the optimal $k_g$ shifts upward when the local baseline is in
play. Finally, the options-framework lens unifies all three: the $1/k_g$
factor is the SMDP discount, the local baseline is the option-level critic,
and the variance bound becomes a special case of the option-Bellman fixed
point's contraction rate.

\bibliographystyle{plainnat}
\bibliography{references}

\end{document}
